\SourcePage{190}{193}
\section{Transitions under a periodic perturbation}

A different type of result is obtained for the probability of transition to
continuous-spectrum states under a periodic perturbation. Suppose that at an
initial time $t=0$ the system is in the $i$th stationary state of the discrete
spectrum. Let the frequency $\omega$ of the periodic perturbation satisfy
\begin{equation}
  \hbar\omega>E_{\min}-E_i^{(0)},
  \tag{42.1}
\end{equation}
where $E_{\min}$ is the energy at which the continuous spectrum begins.

It follows at once from the results of \secref{40} that the principal role is played
by continuous-spectrum states whose energies $E_f$ lie close to the
``resonant'' energy $E_i^{(0)}+\hbar\omega$, that is, states for which
$\omega_{fi}-\omega$ is small. For the same reason, only the first term in the
perturbation matrix element (40.8), whose frequency
$\omega_{fi}-\omega$ is close to zero, need be retained. Substituting this
term in (40.5) and integrating gives
\begin{equation}
  a_{fi}=-\frac{i}{\hbar}\int_0^t V_{fi}(t)\,dt
  =-F_{fi}\frac{\exp[i(\omega_{fi}-\omega)t]-1}
                  {\hbar(\omega_{fi}-\omega)}.
  \tag{42.2}
\end{equation}
\SourcePage{191}{194}
The lower integration limit has been chosen so that $a_{fi}=0$ at $t=0$, in
accordance with the initial condition.

The squared modulus of $a_{fi}$ is
\begin{equation}
  |a_{fi}|^2=|F_{fi}|^2
  \frac{4\sin^2\dfrac{\omega_{fi}-\omega}{2}t}
       {\hbar^2(\omega_{fi}-\omega)^2}.
  \tag{42.3}
\end{equation}
It is readily seen that, at large $t$, the function appearing here may be
represented as proportional to $t$.

To see this, use the formula
\begin{equation}
  \lim_{t\to\infty}\frac{\sin^2\alpha t}{\pi t\alpha^2}=\delta(\alpha).
  \tag{42.4}
\end{equation}
For $\alpha\ne0$ the displayed limit is zero, while for $\alpha=0$,
\[
  \frac{\sin^2\alpha t}{t\alpha^2}=t,
\]
so the limit is infinite. On integrating over $d\alpha$ from $-\infty$ to
$+\infty$ and making the substitution $\alpha t=\xi$, we obtain
\[
  \frac1\pi\int_{-\infty}^{+\infty}
  \frac{\sin^2\alpha t}{t\alpha^2}\,d\alpha
  =\frac1\pi\int_{-\infty}^{+\infty}
  \frac{\sin^2\xi}{\xi^2}\,d\xi=1.
\]
Thus the function on the left-hand side of (42.4) satisfies all the
requirements defining the delta function.

Accordingly, for large $t$ we may write
\[
  |a_{fi}|^2=\frac1{\hbar^2}|F_{fi}|^2\pi t\,
  \delta\left(\frac{\omega_{fi}-\omega}{2}\right),
\]
or, substituting $\hbar\omega_{fi}=E_f-E_i^{(0)}$ and using
$\delta(ax)=\delta(x)/a$,
\[
  |a_{fi}|^2=\frac{2\pi}{\hbar}|F_{fi}|^2
  \delta(E_f-E_i^{(0)}-\hbar\omega)t.
\]

The quantity $|a_{fi}|^2d\nu_f$ is the probability of transition from the
initial state to states in a specified interval $d\nu_f$. At large $t$ it is
proportional to the interval of
\SourcePage{192}{195}
time elapsed since $t=0$. The transition probability $dw_{fi}$ per unit time
is therefore\footnote{It is readily verified that retaining the omitted
second term in (40.8) would give additional expressions which, after division
by $t$, tend to zero as $t\to+\infty$.}
\begin{equation}
  dw_{fi}=\frac{2\pi}{\hbar}|F_{fi}|^2
  \delta(E_f-E_i^{(0)}-\hbar\omega)\,d\nu_f.
  \tag{42.5}
\end{equation}

As expected, it is nonzero only for transitions to states with energy
$E_f=E_i^{(0)}+\hbar\omega$. If the continuous-spectrum energy levels are
nondegenerate, so that $\nu_f$ may be taken to mean the energy alone, the
whole ``interval'' $d\nu_f$ reduces to a single state of energy
$E=E_i^{(0)}+\hbar\omega$, and the probability of transition to this state is
\begin{equation}
  w_{Ei}=\frac{2\pi}{\hbar}|F_{Ei}|^2.
  \tag{42.6}
\end{equation}

It is also instructive to derive (42.5) in another way. Instead of switching
the periodic perturbation on at the definite instant $t=0$, let it grow
slowly from $t=-\infty$ according to the exponential law $e^{\lambda t}$,
where $\lambda$ is positive and is ultimately allowed to tend to zero
(\emph{adiabatic switching}). The initial condition $a_{fi}=0$ is then imposed
at $t=-\infty$. The perturbation matrix element is now
\[
  V_{fi}=F_{fi}e^{i(\omega_{fi}-\omega)t+\lambda t},
\]
and in place of (42.2) we write
\begin{equation}
  a_{fi}=-\frac{i}{\hbar}\int_{-\infty}^{t}V_{fi}(t)\,dt
  =-F_{fi}\frac{\exp[i(\omega_{fi}-\omega)t+\lambda t]}
                  {\hbar(\omega_{fi}-\omega-i\lambda)}.
  \tag{42.7}
\end{equation}
Hence
\[
  |a_{fi}|^2=\frac1{\hbar^2}|F_{fi}|^2
  \frac{e^{2\lambda t}}{(\omega_{fi}-\omega)^2+\lambda^2}.
\]
The transition probability per unit time is determined by
\[
  \frac d{dt}|a_{fi}|^2=2\lambda|a_{fi}|^2.
\]
Now use the formula
\begin{equation}
  \lim_{\lambda\to0}\frac{\lambda}
  {\pi(\alpha^2+\lambda^2)}=\delta(\alpha),
  \tag{42.8}
\end{equation}
\SourcePage{193}{196}
which is valid in the same sense as (42.4). Taking the limit
$\lambda\to0$ gives
\[
  \frac d{dt}|a_{fi}|^2\longrightarrow
  \frac{2\pi}{\hbar^2}|F_{fi}|^2\delta(\omega_{fi}-\omega),
\]
and we return once more to (42.5).
